\chapter{Backend FastAPI}\label{cap:backend}

\section{Estructura del proyecto}

\begin{lstlisting}[caption={Layout del backend}]
fakenews-backend/
|-- main.py            # Puntos de entrada y endpoints REST
|-- billing.py         # Logica de planes, cuotas y Stripe
|-- fever/             # Agente de verificacion
|   |-- agent.py
|   |-- aggregation.py
|   |-- claim_extraction.py
|   |-- inference.py
|   |-- retrieval.py
|   `-- schemas.py
|-- migrations/        # Esquema SQL versionado
|-- tests/             # Pytest
|-- Dockerfile
|-- requirements.txt
`-- requirements-dev.txt
\end{lstlisting}

\section{Endpoints principales}

\begin{itemize}[leftmargin=*,itemsep=0.2em]
  \item \texttt{GET /health} — verificación de estado.
  \item \texttt{POST /predecir/} — clasificación clásica fake/real,
        consume cuota.
  \item \texttt{POST /verify} — agente FEVER (\emph{Super Pro}),
        consume cuota.
  \item \texttt{POST /analyses/save} — persistencia de un análisis en
        el historial del usuario.
  \item \texttt{POST /account/delete} — borrado de cuenta y datos
        asociados.
  \item Endpoints de \emph{billing}: creación de \emph{checkout},
        consulta de plan, \emph{webhook} de Stripe.
\end{itemize}

\section{Validación e inyección de dependencias}

Cada endpoint declara los \emph{schemas} de entrada y salida con
Pydantic, lo que garantiza validación automática y generación de
OpenAPI. La autenticación se resuelve con una dependencia que extrae
el JWT del encabezado \texttt{Authorization}, lo verifica contra las
claves públicas de Supabase y devuelve el \emph{usuario actual}.

\section{Gestión de cuotas}

La función \texttt{\_consume\_verification\_quota} encapsula:
\begin{enumerate}[leftmargin=*,itemsep=0.2em]
  \item Lectura del plan vigente del usuario (tabla
        \texttt{billing.subscriptions}).
  \item Cómputo de uso del periodo (tabla \texttt{usage}).
  \item Comprobación de cuota; respuesta HTTP 402 / 403 si excede.
  \item Registro atómico de un nuevo evento de uso.
\end{enumerate}

\section{Pruebas}

La carpeta \path{tests/} contiene 25 tests (Pytest) cubriendo
\emph{billing}, \emph{smoke} y agente. Se ejecutan en cada
\emph{commit} mediante \texttt{pytest -q}.

\section{Empaquetado}

\path{Dockerfile} parte de una imagen Python 3.12 \emph{slim},
instala dependencias en una capa cacheada y expone el puerto 8000 vía
\texttt{uvicorn}.
